% !TeX document-id = {1c0b4298-276a-4038-8e08-c4a1f4846da7}
% !TeX TXS-program:bibliography = txs:///biber
\documentclass[12pt,a4paper]{article}
\usepackage[T1]{fontenc}
\usepackage{graphicx}
\usepackage{mathtools}
\usepackage{amssymb}
\usepackage{amsthm}
\usepackage{thmtools}
\usepackage{xcolor}
\usepackage{nameref}
\usepackage{hyperref}
\usepackage{color}
\usepackage{float}
\newtheorem{theorem}{Teorema}
%\usepackage[margin=2.5cm]{geometry}
\usepackage[brazilian]{babel}
\usepackage[
backend=biber,
style=authoryear,
natbib=true
]{biblatex}
\addbibresource{../bibliography.bib}
\newtheorem{definition}{Definição}
\newtheorem{lemma}{Lema}
\newtheorem{proposition}{Proposição}
\title{Demonstração do teorema de Gauss-Markov}
\author{ Professor Luís Antonio Fantozzi Alvarez}
\date{15 de abril de 2026}
\begin{document}
	\maketitle
	
	Nestas notas, faremos a demonstração do teorema de Gauss-Markov visto em aula, cujo enunciado reproduzimos a seguir:
	\begin{theorem}[Gauss-Markov]
		Sob H1-H3, o estimador de MQO é o estimador {\color{red}linear} não viciado para $\delta(F)$ de variância uniformemente mínima na classe $\mathcal{F}(\boldsymbol{X})$, i.e. para qualquer outro estimador  {\color{red}linear} $\psi$ tal que $\mathbb{E}_{F}[\psi|\boldsymbol{X}] = \delta(F)$ para todo $F \in \mathcal{F}(\boldsymbol{X})$, temos que, para todo $F \in \mathcal{F}(\boldsymbol{X})$:
		
		$$\mathbb{V}_{F}(\psi|\boldsymbol{X}) - \mathbb{V}_{F}(\hat{b}|\boldsymbol{X})\, \quad \text{é positiva semidefinida} $$
	\end{theorem}
	
\begin{proof}
Dividimos a demonstração em quatro etapas.
\paragraph{Etapa 1} Observe que $\{\delta(F) : F \in \mathcal{F}(\boldsymbol(X))\} = \mathbb{R}^k$, pois, para todo $\delta \in \mathbb{R}^k$, conseguimos definir um modelo dado por $\boldsymbol{y}=\boldsymbol{X}\delta + \boldsymbol{\epsilon}$, onde $\boldsymbol{\epsilon} \sim \mathcal{N}(\boldsymbol{0}_{n\times 1},\mathbb{I}_{n\times n})$, com $\boldsymbol{\epsilon}$ independente de $\boldsymbol{X}$, e a distribuição $P_{\boldsymbol{y}|\boldsymbol{X}}$ resultante pertencerá a $\mathcal{F}(\boldsymbol{X})$, com $\delta(P_{\boldsymbol{y}|\boldsymbol{X}})=\delta$.
\paragraph{Etapa 2} Considere um estimador linear $A(\boldsymbol{X})\boldsymbol{y}$ não vixiado para $\delta(F)$ em $\mathcal{F}(\boldsymbol{X})$. Se esse estimador é não viciado, então $\mathbb{E}_F[A(\boldsymbol{X})\boldsymbol{y}|\boldsymbol{X}]=A(\boldsymbol{X})\boldsymbol{X}\delta(F) = \delta(F)$ para toda $F \in \mathcal{F}(\boldsymbol{X})$. Fixe $j \in \{1,\ldots, k\}$. Tomando $F$ tal que $\delta(F) = e_j$, onde $e_j$ é o vetor $k\times 1$ com $1$ na $j$-ésima entrada e $0$ nas demais, concluímos então que a $j$-ésima coluna de $A(\boldsymbol{X})\boldsymbol{X}$ é igual a $e_j$. Como $j$ foi escolhido arbitrariamente, concluímos que $A(\boldsymbol{X})\boldsymbol{X} = \begin{bmatrix}e_1 & e_2& \ldots & e_k\end{bmatrix} = \mathbb{I}_{k \times k}$ .
\paragraph{Etapa 3} Ainda considerando o estimador da etapa anterior, fixe $F \in \mathcal{F}(\boldsymbol{X})$. Observamos que: 

\begin{equation*}
\begin{aligned}
\operatorname{cov}_F(A(\boldsymbol{X})\boldsymbol{y}-\hat{b},\hat{b}|\boldsymbol{X}) =\\ \mathbb{E}_F[(A(\boldsymbol{X})(\boldsymbol{y}-\mathbb{E}_F[\boldsymbol{y}|\boldsymbol{X}])-(\hat{b}-\mathbb{E}_F[\hat{b}|\boldsymbol{X}]))(\hat{b}-\mathbb{E}_F[\hat{b}|\boldsymbol{X}])'|\boldsymbol{X}] = \\
\mathbb{E}_F[(A(\boldsymbol{X}) - (\boldsymbol{X}'\boldsymbol{X})^{-1}\boldsymbol{X}')\boldsymbol{\epsilon}\boldsymbol{\epsilon}'\boldsymbol{X} (\boldsymbol{X}'\boldsymbol{X})^{-1}|\boldsymbol{X}] = \\
\sigma^2(\underbrace{A(\boldsymbol{X})\boldsymbol{X}}_{=\mathbb{I}_{k\times k}} (\boldsymbol{X}'\boldsymbol{X})^{-1} -       \underbrace{(\boldsymbol{X}'\boldsymbol{X})^{-1}\boldsymbol{X}' \boldsymbol{X}}_{=\mathbb{I}_{k\times k}} (\boldsymbol{X}'\boldsymbol{X})^{-1}) = 0_{k\times k}
\end{aligned}
\end{equation*}

Como a covariância acima é zero, concluímos que: $\operatorname{cov}_F(A(\boldsymbol{X})\boldsymbol{y},\hat{b}|\boldsymbol{X})=\mathbb{V}_F[\hat{b}|\boldsymbol{X}]$.
\paragraph{Etapa 4:} Finalmente notamos que, usando a Etapa 3:

\begin{equation*}
\begin{aligned}\mathbb{V}_F(A(\boldsymbol{X})\boldsymbol{y} -\hat{b}|\boldsymbol{X} ) =\\ \mathbb{V}_F(A(\boldsymbol{X})\boldsymbol{y}|\boldsymbol{X} ) +\mathbb{V}_F(\hat{b}|\boldsymbol{X} ) -\operatorname{cov}_F(A(\boldsymbol{X})\boldsymbol{y},\hat{b}|\boldsymbol{X}) - \operatorname{cov}_F(\hat{b},A(\boldsymbol{X})\boldsymbol{y}|\boldsymbol{X}) =\\  \mathbb{V}_F(A(\boldsymbol{X})\boldsymbol{y}|\boldsymbol{X} )-\mathbb{V}_F(\hat{b}|\boldsymbol{X} )\, ,\end{aligned}
\end{equation*}
e, como $\mathbb{V}_F(A(\boldsymbol{X})\boldsymbol{y} -\hat{b}|\boldsymbol{X} ) $ é uma matriz positiva semidefinida, obtemos a conclusão desejada.
\end{proof}
\end{document}